% !Mode:: "TeX:UTF-8"

\documentclass[fontset=windows,type=master,campus=shenzhen, chapterbold=false, tocblank=false]{hithesisbook}


\usepackage{hithesis}
\newcommand{\cndash}{\rule{0.2em}{0pt}\rule[0.35em]{1.6em}{0.05em}\rule{0.2em}{0pt}}


\graphicspath{{figures/}}

% ===== 算法环境说明 =====
% 模板 hithesis.sty 已加载 algorithm2e 并完成全部配置：
%   algoruled       — 三线格式：上1.5pt细线 + 标题 + 下1.5pt粗线
%   linesnumbered   — 自动行号
%   algochapter     — 按章编号（算法3-1、算法3-2...）
%   \algorithmcfname=算法   — 中文"算法"前缀
%   标题：\wuhao（五号），算法体：继承正文 \normalsize（小四）
% 使用方式：\begin{algorithm}[htbp] \caption{名称} \label{alg:xxx} ... \end{algorithm}
\DontPrintSemicolon

\begin{document}
\frontmatter
\input{front/cover} % 封面
\makecover
% \input{front/denotation}%物理量名称表,符合规范为主，有要求添加
\tableofcontents %目录
\mainmatter
% !Mode:: "TeX:UTF-8"

\chapter{绪论}[Introduction]

\section{研究背景与意义}

本文使用哈尔滨工业大学深（深圳）深圳学位论文模板~\cite{hithesis2017} 进行排版。正常些据可以


\section{国内外研究现状}

\section{本文主要研究内容}

\section{论文组织结构}
   % 绪论
% !Mode:: "TeX:UTF-8"

\chapter{相关理论与技术基础}[Related Theory and Technical Foundation]

% TODO: 填写相关技术基础内容

\section{STL模型与增材制造}

\section{拓扑数据结构}

\section{哈希索引与空间查询}

\section{本章小结}
   % 悬垂区域识别与建造方向优化
% !Mode:: "TeX:UTF-8"

\chapter{三维网格切片算法}[3D Mesh Slicing Algorithm]
   % 三维网格切片算法
% !Mode:: "TeX:UTF-8"

\chapter{实验与结果分析}[Experiments and Results Analysis]
   % 实验与结果分析
\backmatter
% !Mode:: "TeX:UTF-8" 
\begin{conclusions}

学位论文的结论作为论文正文的最后一章单独排写，但不加章标题序号。

结论应是作者在学位论文研究过程中所取得的创新性成果的概要总结，不能与摘要混为一谈。

\end{conclusions}
   % 结论
\bibliographystyle{gbt7714-numerical}




\bibliography{reference} % 参考文献
% \begin{appendix}%附录
% \input{back/appA.tex}
% \end{appendix}
\include{back/publications}    % 所发文章
% \include{back/ceindex}    % 索引, 根据自己的情况添加或者不添加，选择自动添加或者手工添加。
\authorization %授权
% !Mode:: "TeX:UTF-8"
\begin{acknowledgements}

衷心感谢哈尔滨工业大学（深圳）各位领导为毕业设计模板撰写提供宝贵机会；特别感谢教务处老师、王鸿鹏教授和夏文教授对本模版给予的专业建议与支持；同时感谢
hitszhesis
和
hithesis
项目提供的优质工具与参考资料。


\end{acknowledgements}
 %致谢
\end{document}
% Local Variables:
% TeX-engine: xetex
% End:
